\section{Positivity Off the Critical Plane}
\label{sec:positivity}

The heart of our proof is showing that the operator $I - A(s)$ is positive definite for $\text{Re}(s) > n+1$. This follows from Recognition Science's Positive Cost axiom (A3), which states that all recognition events carry non-negative cost.

\subsection{The Norm Bound}

\begin{lemma}[Operator Norm]
\label{lem:norm-bound}
For $s = \sigma + it$ with $\sigma = \text{Re}(s)$, we have
\[
\|A(s)\| = 2^{-((\sigma-(n+1))+\varepsilon)}.
\]
\end{lemma}

\begin{proof}
Since $A(s)$ is diagonal, its operator norm equals the supremum of the moduli of its eigenvalues:
\[
\|A(s)\| = \sup_{q \in \mathcal{P}} |q^{-(s-(n+1)+\varepsilon)}| = \sup_{q \in \mathcal{P}} q^{-((\sigma-(n+1))+\varepsilon)}.
\]
The supremum is achieved at $q = 2$ (the smallest prime power), giving $\|A(s)\| = 2^{-((\sigma-(n+1))+\varepsilon)}$.
\end{proof}

\subsection{Positivity from Recognition Cost}

The key insight is that the operator $I - A(s)$ represents the ledger cost of maintaining Hodge patterns at energy scale $s$.

\begin{theorem}[Positivity Off Critical Plane]
\label{thm:positivity}
For $\text{Re}(s) > n+1$, the operator $I - A(s)$ is positive definite.
\end{theorem}

\begin{proof}
We need to show that $\langle (I - A(s))f, f \rangle > 0$ for all non-zero $f \in \mathcal{H}$.

For $\sigma = \text{Re}(s) > n+1$, we have $(\sigma-(n+1))+\varepsilon > \varepsilon = \varphi - 1 > 0$.

By Lemma \ref{lem:norm-bound}, $\|A(s)\| = 2^{-((\sigma-(n+1))+\varepsilon)} < 2^{-\varepsilon} < 1$.

For any $f \in \mathcal{H}$ with $\|f\| = 1$:
\[
\langle (I - A(s))f, f \rangle = \langle f, f \rangle - \langle A(s)f, f \rangle = 1 - \langle A(s)f, f \rangle.
\]

Since $A(s)$ has positive eigenvalues when $s$ is real (Lemma \ref{lem:operator-props}(iii)), and by continuity for complex $s$ with $\text{Re}(s) > n+1$:
\[
|\langle A(s)f, f \rangle| \leq \|A(s)\| \cdot \|f\|^2 = \|A(s)\| < 1.
\]

Therefore, $\langle (I - A(s))f, f \rangle \geq 1 - \|A(s)\| > 0$.
\end{proof}

\subsection{Spectral Analysis}

Since $A(s)$ is diagonal, we can analyze the spectrum explicitly.

\begin{lemma}[Spectral Gap]
\label{lem:spectral-gap}
For $\text{Re}(s) > n+1$, the operator $I - A(s)$ has spectral gap
\[
\text{gap}(I - A(s)) = 1 - 2^{-((\text{Re}(s)-(n+1))+\varepsilon)} > 0.
\]
\end{lemma}

\begin{proof}
The eigenvalues of $I - A(s)$ are $\{1 - q^{-(s-(n+1)+\varepsilon)} : q \in \mathcal{P}\}$.

For $\text{Re}(s) > n+1$:
- The largest eigenvalue of $A(s)$ is $2^{-((\text{Re}(s)-(n+1))+\varepsilon)}$ (at $q = 2$)
- The smallest eigenvalue of $I - A(s)$ is $1 - 2^{-((\text{Re}(s)-(n+1))+\varepsilon)}$
- All other eigenvalues of $I - A(s)$ are larger

The spectral gap is therefore $1 - 2^{-((\text{Re}(s)-(n+1))+\varepsilon)} > 0$.
\end{proof}

\subsection{Connection to Recognition Science}

In the Recognition Science framework, the positivity of $I - A(s)$ has a deep interpretation:

\begin{proposition}[Recognition Cost Interpretation]
The quadratic form $\langle (I - A(s))f, f \rangle$ represents the total recognition cost of maintaining the Hodge pattern $f$ at energy scale $s$. Positivity for $\text{Re}(s) > n+1$ means that all non-trivial patterns have positive maintenance cost above the critical energy.
\end{proposition}

This aligns with Recognition Science Axiom A3: all physical processes (including pattern maintenance) require positive energy cost. The critical line $\text{Re}(s) = n+1$ represents the boundary where patterns can exist with zero cost—these are precisely the algebraic cycles.

\subsection{Determinant Non-Vanishing}

As an immediate consequence of positivity:

\begin{corollary}[Determinant Non-Zero]
\label{cor:det-nonzero}
For $\text{Re}(s) > n+1$, we have $\det_2(I - A(s)) > 0$.
\end{corollary}

\begin{proof}
A positive definite operator on a Hilbert space has strictly positive determinant. Since $I - A(s)$ is positive definite for $\text{Re}(s) > n+1$ by Theorem \ref{thm:positivity}, and $A(s)$ is Hilbert-Schmidt in this region (Lemma \ref{lem:operator-props}(ii)), the regularized determinant $\det_2(I - A(s))$ is well-defined and strictly positive.
\end{proof}

This non-vanishing will be crucial in the next section when we analyze the zeros of $\zeta_{\text{Hdg}}(s)$. 